\documentclass[11pt]{article}

\usepackage[margin=1in]{geometry}
\usepackage{amsmath,amssymb,amsthm}
\usepackage{mathtools}
\usepackage{mathrsfs}
\usepackage{enumitem}
\usepackage{hyperref}
\usepackage{xcolor}

\newtheorem{theorem}{Theorem}
\newtheorem{lemma}[theorem]{Lemma}
\newtheorem{proposition}[theorem]{Proposition}
\newtheorem{corollary}[theorem]{Corollary}
\theoremstyle{definition}
\newtheorem{definition}[theorem]{Definition}
\newtheorem{remark}[theorem]{Remark}

\newcommand{\R}{\mathbb{R}}
\newcommand{\C}{\mathbb{C}}
\newcommand{\Z}{\mathbb{Z}}
\newcommand{\KG}{\mathcal{K}_G}
\newcommand{\Heta}{H_{\eta}}
\newcommand{\HzeroP}{H_0'}
\newcommand{\HzeroPP}{H_0''}
\newcommand{\HetaP}{H_{\eta}'}
\newcommand{\HetaPP}{H_{\eta}''}
\newcommand{\Phizz}{\Phi}
\newcommand{\Fmap}{F}
\newcommand{\Finvmap}{F^{-1}}
\newcommand{\Om}{\Omega}
\newcommand{\dOm}{\partial\Omega}
\newcommand{\sgn}{\operatorname{sgn}}
\newcommand{\norm}[1]{\left\lVert#1\right\rVert}
\newcommand{\abs}[1]{\left\lvert#1\right\rvert}
\newcommand{\abslb}[1]{\underline{|#1|}}
\newcommand{\absub}[1]{\overline{|#1|}}

\title{Novel mathematical insights from the proof-refinement loop \\
       (Lemma~5.4 of Braverman~et~al., \emph{The Grothendieck constant is strictly smaller than Krivine's bound})}
\date{}

\begin{document}
\maketitle

\noindent\textit{Appendix to the main paper.}
This appendix collects the genuine mathematical content discovered while driving the
constant $C_r$ in the rigorous certificate of Lemma~5.4 down through several orders of
magnitude. We follow the convention that $\KG$ denotes the (real) Grothendieck constant.
The Krivine upper bound is $\KG \le \pi/(2\ln(1+\sqrt 2))\approx 1.7822$;
each refinement of the inequality below produces a strictly smaller certified upper
bound for $\KG$. Throughout, $\eta=1/100$ and $r=23/25$ unless stated otherwise, and
$\Om(\delta,B)\subset\C$ denotes the rectangle $\{w:|\Re w|\le 1-\delta,\ |\Im w|\le B\}$
on which the implicit-function family is analytic.

\bigskip

%-------------------------------------------------------------------------
\section{The improvement-over-Krivine quantity}\label{sec:KG-quantity}

In the Braverman--Guruswami--Lee--Naor program, the certified upper bound on $\KG$
takes the form
\begin{equation}\label{eq:KG-bound}
   \KG \;\le\; \frac{1}{c_{f,g}}, \qquad
   c_{f,g} = \tfrac{2}{\pi}\,\gamma, \qquad
   \gamma = \gamma_0 + \Delta\gamma,
\end{equation}
where $\gamma_0=\operatorname{atanh}(1/\sqrt 2)=\ln(1+\sqrt 2)$ is the value the inequality
collapses to in the absence of the Lemma~5.4 refinement. Setting $\Delta\gamma = 0$ recovers
the classical Krivine bound $\pi/(2\ln(1+\sqrt 2))$. The whole program is the
construction of a strictly positive, rigorously certified $\Delta\gamma>0$. The natural
\emph{quantity of merit} is therefore the gap between Krivine's constant and the new
upper bound,
\begin{equation}\label{eq:KGimp}
   \KG\text{-improvement}
   \;\coloneqq\;
   \frac{\pi}{2\ln(1+\sqrt 2)} \;-\; \frac{1}{c_{f,g}}
   \;=\;
   \frac{\Delta\KG}{c_{\text{base}}\,(c_{\text{base}}+\Delta\KG)},
\end{equation}
with $c_{\text{base}}=(2/\pi)\gamma_0$ and $\Delta\KG=(2/\pi)\Delta\gamma$. Two
properties of \eqref{eq:KGimp} are used systematically in the audits and figures.

\paragraph{Inversion stability.}
Because $c_{\text{base}}=(2/\pi)\ln(1+\sqrt 2) \approx 0.5611$ is bounded away from zero,
the radius of the interval $1/c_{f,g}$ in arb (interval-ball) arithmetic is comparable to
the radius of $c_{f,g}$ itself: a relative widening by the factor
$1/c_{\text{base}}^2\approx 3.18$. Subtraction with the constant $\pi/(2\gamma_0)$ does
\emph{not} suffer cancellation in arb because the radius is bounded above by the
upstream stored radius $\sim 10^{-41}$, well below the actual gap $\Delta\KG/c_{\text{base}}^2$
which lies between $10^{-32}$ and $10^{-17}$ across iterates. This guarantees that the
plotted improvement values are not artifacts of floating-point cancellation.

\paragraph{Linearisation.}
For $\Delta\KG\ll c_{\text{base}}$,
\eqref{eq:KGimp} expands as
\(
   \KG\text{-improvement}
   = (1/c_{\text{base}}^2)\,\Delta\KG + O(\Delta\KG^2)
   = \bigl(\pi/(2\gamma_0)\bigr)^2 \Delta\KG + O(\Delta\KG^2),
\)
so the curve $\KG$-improvement vs.\ iteration is, to leading order, just the Krivine
constant squared times the $\Delta\KG$ trace. The plots in
\href{run:plots/kg_improvement.pdf}{\texttt{plots/kg\_improvement.pdf}} display the
exact form \eqref{eq:KGimp} computed with arb at precision~256.

\bigskip

%-------------------------------------------------------------------------
\section{The auxiliary $z$-identity for $\norm{u_{zz}}$}\label{sec:aux-z-identity}

The original argument bounds $\norm{u_{zz}}$ via a recentred Cauchy estimate on the
analyticity disk $\{|z-z_0|<r\}$:
\(
   \norm{u_{zz}}\le 8R/(1-r)^2.
\)
The new derivation comes from differentiating the implicit-function identity
$\Fmap(\eta, p, \Finvmap(p,z)) = z$ twice in $z$ and exploiting the algebraic relation
\begin{equation}\label{eq:auxz}
   D\cdot u_{zz}
   \;+\;
   DD\cdot h_z^2
   \;=\; 0,
   \qquad
   D = \partial_w \Fmap,\quad
   DD = \partial_w^2 \Fmap,\quad
   h_z = \partial_z\Finvmap = 1/D.
\end{equation}
Solving \eqref{eq:auxz} gives the \emph{exact pointwise identity}
\begin{equation}\label{eq:uzz-id}
   u_{zz}
   \;=\;
   -\,\frac{DD\,h_z^2}{D}
   \;=\;
   -\,\frac{DD}{D^3}.
\end{equation}
The auxiliary $z$-identity \eqref{eq:auxz} is, up to sign, a re-statement of
$\partial_z(D h_z) = 0$ along the inverse curve.
Identity~\eqref{eq:uzz-id} replaces the worst-case Cauchy step
$8R/(1-r)^2 \approx 1393$ by a direct sup
\(
   \norm{u_{zz}}\le \norm{DD/D^3}_{\Om}
\)
which our truncation-free chunk certifier evaluates to a value $\approx 1.55$,
a $\sim 900\times$ tightening of the term~that dominated the V10 baseline.
Crucially this is identity (not inequality) at every $(p,w)$, so the slack is purely a
function of how tightly we can sup over a given $(p,w)$-cell.

\bigskip

%-------------------------------------------------------------------------
\section{Joint sup-bounds vs.\ triangle products}\label{sec:joint-vs-triangle}

A recurring leverage point is that the Lemma~5.4 derivation produces inequalities of
the form
\begin{equation}\label{eq:triangle}
   \sup_{(p,w)} \norm{ A(p,w)\,B(p,w) /\ C(p,w)^k }
   \;\le\;
   \Bigl(\sup A\Bigr)\,\Bigl(\sup B\Bigr) / \Bigl(\inf C\Bigr)^k,
\end{equation}
and the right-hand side is what the original proof bounds by \emph{three independent}
suprema. \eqref{eq:triangle} is sharp only when the three suprema all bind at the same
$(p,w)$. For our analytic family they bind at \emph{different} corners, and the
right-hand side over-estimates the joint sup by a factor that turns out to be
$5$--$300\times$ depending on the cell.

\paragraph{Term~1 (sprint~019).}
The triangle product
\(
M_\phi\,M_{\phi'}\,(1/D)^2
\)
on the V13 oracle constellation evaluates to $2.09\times 10^{-3}$, while the joint
sup
\(
J_1 \;=\; \sup_{(p,w)} \norm{\Phizz(w)\,\Phizz'(w)\,/\,D(p,w)^2}
\)
certifies to $J_1 \le 3\times 10^{-4}$ on the same rectangle---a $\sim 7\times$
reduction. Mathematically this happens because $\Phi(w)$ and $\Phi'(w)$ are
\emph{$\eta$-correlated}: each is an $O(\eta^2)$ deviation from
$H_0(w)$ resp.\ $H_0'(w)$, and the deviations have aligned phase, so $|\Phi\,\Phi'|$ at
the binding $w$ is much smaller than the product of the worst-case magnitudes.

\paragraph{Term~2 (sprint~020).}
The same mechanism applied to the second term:
\(
   \tfrac12 u_{zz}\,M_\phi^2 = \tfrac12\,\norm{DD/D^3}\,M_\phi^2 \approx 9.24\times 10^{-4}
\)
becomes
\(
   J_2 \;=\;\tfrac12\sup_{(p,w)} \norm{ DD(p,w)\,\Phi(w)^2 / D(p,w)^3 } \approx 7.7\times 10^{-5}
\),
a $\sim 6\times$ reduction. The \emph{same} aux-$z$ identity $u_{zz} = -DD/D^3$
explains why $u_{zz}\,\Phi^2 = -DD\,\Phi^2/D^3$ is the natural single-cell
quantity to certify.

\bigskip

%-------------------------------------------------------------------------
\section{The form-A identity and merged-numerator interference}\label{sec:formA-merged}

\paragraph{Form-A identity.}
Differentiating $\Fmap(\eta,p,\Finvmap(p,z))=z$ twice in $p$ at fixed $z$ yields
\begin{equation}\label{eq:formA}
   D\,u_{pp} \;+\; 2\,\Phi'\,u_p \;+\; DD\,u_p^2 \;=\; 0,
\end{equation}
where $u_p = \partial_p \Finvmap$ satisfies, by a separate differentiation,
$u_p = -\Phi/D$ with $\Phi = \Heta - H_0$. Substituting $u_p = -\Phi/D$ into
\eqref{eq:formA} and solving gives the \emph{exact pointwise identity}
\begin{equation}\label{eq:upp-merged}
   u_{pp}
   \;=\;
   \frac{2\,\Phi(w)\,\Phi'(w)\,D(p,w)\;-\;DD(p,w)\,\Phi(w)^2}{D(p,w)^3}.
\end{equation}
This is the \emph{merged numerator}. The two terms in the numerator have opposite
signs and partially cancel.

\paragraph{Triangle vs.\ merged.}
A naive triangle bound on $|u_{pp}|$ would split \eqref{eq:upp-merged} into
\(
   |u_{pp}|
   \;\le\;
   2\,\norm{\Phi\Phi'/D^2} + \norm{DD\,\Phi^2/D^3}
   \;=\;
   2 J_1 + J_2,
\)
discarding the destructive interference between the two summands. The merged sup
\(
   J_{\mathrm{merged}}
   \;=\;
   \sup_{(p,w)} \frac{ |2\Phi\Phi'D - DD\,\Phi^2| }{ |D|^3 }
\)
is the genuine pointwise sup of $|u_{pp}|$, and $J_{\mathrm{merged}} \le 2J_1+J_2$
with equality if the two summands are everywhere phase-aligned.

\paragraph{Resolution-dependent crossover.}
A surprising phenomenon emerges when computing $J_{\mathrm{merged}}$ rigorously by
interval arithmetic on cells of $\dOm \times [0,1]$ at resolution~$n$:
\[
\begin{array}{c|cccccccc}
   n & 1{,}024 & 2{,}048 & 4{,}096 & 8{,}192 & 16{,}384 & 32{,}768 & 65{,}536 & 131{,}072 \\\hline
   J_{\mathrm{merged}}/(2J_1+J_2) & 1.99 & 1.92 & 1.87 & 1.78 & 1.61 & 1.40 & 1.20 & 1.034
\end{array}
\]
At low resolution the certified bound on the \emph{difference} $|2\Phi\Phi'D-DD\Phi^2|$
is dominated by the per-cell interval slack of subtracting two large quantities of
opposite sign---the bound is \emph{worse} than the triangle. The cancellation is real
but invisible at coarse scale. Around $n\approx 2^{17}\!-\!2^{18}$ the per-cell slack
shrinks below the cancellation gain and the merged bound crosses below the triangle.
At $n=262{,}144$ the certified merged bound is $\approx 0.86\times$ the triangle,
yielding a structural $\sim 14\%$ reduction that the triangle decomposition cannot
capture at any resolution.

This crossover is the cleanest example we have observed of a rigorous bound that is
strictly tighter only above a quantitative resolution threshold, where ``resolution''
refers to the discretisation of the sup.

\bigskip

%-------------------------------------------------------------------------
\section{The $K_{\max}$ truncation regression and a truncation-free pipeline}\label{sec:trunc-free}

The semi-analytic step in the original numerical pipeline writes
\begin{equation}\label{eq:phi-series}
   \Heta(w)
   \;=\;
   \arcsin(w) \;-\; \frac{\phi\bigl(\arcsin(w)\bigr)}{\sqrt{1-w^2}},
   \qquad
   \phi(z) \;=\; \sum_{k\ge 0} c_k\,z^k,
\end{equation}
and evaluates $\phi$ by truncating the series at $K_{\max}$. The $b$-series file at our
parameters has $250$ coefficients with empirical pseudo-rate $|b_k|^{1/k}\approx 0.964$
near $k=249$; a truncation at $K_{\max}=80$ leaves a tail that is \emph{not bounded
rigorously} for the magnitudes that arise in $\Phi=\Heta-H_0$ at $\eta=1/100$
($\Phi$ itself is $O(\eta^2)\sim 10^{-4}$, whereas the pseudo-rigorous tail is plausibly
$O(10^{-8})$ but unverified). Interval arithmetic over a truncated polynomial returns
a rigorous bound on the wrong function.

\subsection{Integral representation via Price's theorem}

The replacement is to evaluate $\Heta, \HetaP, \HetaPP$ \emph{without} $\phi$.
Following Price's theorem, the analytic continuation of $\Heta$ to $w\in\C$ admits
the integral representation
\begin{equation}\label{eq:price}
   \Heta(w)
   \;=\;
   \frac{1}{2\sigma_z}\,\iint_{\R^2}
       K_z\bigl(P(x), Q(y)\bigr)\,
       e^{-S_z[P,Q](x,y)}\;dx\,dy,
\end{equation}
where $K_z(P,Q)=\sgn(x-\eta P(x))\sgn(y+\eta Q(y))$ for the Hermite-physicist
sign-pair we use, $S_z$ is the action quadratic form parameterised by $z\in\Om$, and
$\sigma_z$ is the corresponding normalising determinant. The first and second
$w$-derivatives are obtained by differentiating \eqref{eq:price} under the integral,
producing analogous formulas for $\HetaP$ and $\HetaPP$ that involve $P', Q'$ but no
power series in $w$.

The integrals \eqref{eq:price} are evaluated via tensor-product Gauss--Legendre at the
chunk midpoint, with adaptive Bernstein-ellipse error per panel; the value is then a
sharp rigorous evaluation of $\Heta(w)$ (and its derivatives) at that midpoint.

\subsection{Real-positive majorants for chunk bounds}

To bound $\sup_{w\in Z}|\Heta(w)|$ over a chunk $Z\subset\Om$, we replace
\eqref{eq:price} with its real-positive majorant
\begin{equation}\label{eq:majorant}
   |\Heta(w)|
   \;\le\;
   \frac{1}{2\absub{\sigma_z}}\,\iint_{\R^2}
       \bigl|K_z(P,Q)\bigr|\;
       \bigl|e^{-S_z}\bigr|\;dx\,dy
   \;\le\;
   \frac{1}{2\abslb{\sigma_z}}\,\iint_{\R^2}
       \widehat E_Z(x,y)\,dx\,dy,
\end{equation}
where $\widehat E_Z(x,y)$ is a real, non-negative Gaussian-shaped majorant of
$|K_z\,e^{-S_z}|$ valid uniformly for $w\in Z$. The right-hand side reduces to a finite
sum of closed-form Gaussian moments
\(
   \int x^{2j} e^{-\alpha x^2}\,dx
\)
that arb evaluates rigorously. The chunk-bound for $\HetaP, \HetaPP, \Phi, \Phi'$ are
analogous, with a key extra ingredient described below.

\subsection{The $\eta^2$-tight $\Phi$-difference bound}

Bounding $\Phi(w)=\Heta(w)-H_0(w)$ by a triangle inequality
$|\Phi|\le |\Heta|+|H_0|\approx 2.4$ throws away the $O(\eta^2)$ scaling of $\Phi$
and produces a $C_r$ that is six orders of magnitude worse than even the heuristic
prior baseline. The remedy is to take the difference \emph{inside} the integral:
\begin{equation}\label{eq:phi-diff}
   \Phi(w)
   \;=\;
   \frac{1}{2\sigma_z}\,\iint_{\R^2}
       \bigl[ K_z(P,Q) - K_z(0,0) \bigr]\;e^{-S_z[P,Q]}\,dx\,dy
       \;+\;\text{(small) Gaussian-action correction},
\end{equation}
and majorise the kernel difference $K_z(P,Q)-K_z(0,0)$ by a real-positive Taylor
expansion in $(P,Q)$ around $(0,0)$. Because $P,Q$ are themselves $O(\eta)$, the
chunk-bound on $\Phi$ scales as $O(\eta^2)$, recovering the cancellation that the
triangle bound discarded.

\subsection{The $\lambda$-decay parallelogram fix}\label{sec:lambda-fix}

A subtle bug in an earlier version of the $\HetaPP$ chunk certifier used the bound
$\lambda \le 2\min(\alpha,\beta)$, where $\alpha,\beta$ are the principal-axis decay
rates of the Gaussian factor. The correct bound is
\begin{equation}\label{eq:lambda-fix}
   \lambda \;\le\; \min(\alpha,\beta),
\end{equation}
which is the parallelogram identity
$(x-y)^2 + (x+y)^2 = 2(x^2+y^2)$ in disguise: the joint quadratic form on the
\emph{difference} of two Gaussians has decay $\min(\alpha,\beta)$, not
$2\min(\alpha,\beta)$. The factor-of-two error inflated the certified Gaussian moments
by $2^{j}$ at the $2j$-th moment and was the main quantitative gap between the
truncation-free pipeline and the heuristic baseline before the fix.

\subsection{Anisotropic Gaussian moments}\label{sec:aniso}

Every chunk-bound integral collapses to a sum of moments of an anisotropic Gaussian
$e^{-(\alpha x^2 + \beta y^2 + 2\gamma xy)}$. Bounding each moment by its isotropic
majorant (replace $\alpha,\beta$ by $\min(\alpha,\beta)$) is sound but loose. We use
\(
   \text{moment}
   \;\le\;
   \min\bigl(\text{isotropic-majorant},\ \text{exact-anisotropic}\bigr),
\)
where the exact-anisotropic value is computed by diagonalising the quadratic form via
an explicit rotation. The minimum is provably an upper bound on the moment and is
tighter than either alone on cells where $\alpha,\beta$ differ substantially.

\bigskip

%-------------------------------------------------------------------------
\section{The denominator-fix for $D(p,w)$}\label{sec:denom-fix}

The joint quotients $\Phi\Phi'/D^2$ and $DD\,\Phi^2/D^3$ require a uniform lower
bound $|D(p,w)|\ge \delta_D > 0$ on each cell. Computing $|D|$ directly as a joint
$(p,w)$ interval inflates by orders of magnitude because $D$ is a difference
$D = \HzeroP + p(\HetaP-\HzeroP)$ where the two summands have similar magnitude.
The fix exploits the affine structure in $p$:
\begin{equation}\label{eq:D-lb}
   |D(p,w)|
   \;\ge\;
   |\HzeroP(w)|
   \;-\;
   |p|\cdot|\HetaP(w) - \HzeroP(w)|
   \;\ge\;
   \abslb{\HzeroP(w)}\;-\;p_{\mathrm{hi}}\cdot M_{\Phi'},
\end{equation}
where $M_{\Phi'}$ is the chunk-certified bound on $|\Phi'|$ and $\abslb{\HzeroP(w)}$ is
the closed-form interval lower bound $1/|1-w^2|^{1/2}$. The right-hand side of
\eqref{eq:D-lb} is computed once per $w$-chunk (independent of $p$) and is positive
provided the chunk is bounded away from the branch points $w=\pm 1$. This avoids
multiplying a small $p$-interval by a large $\HetaP-\HzeroP$ in the joint $(p,w)$
representation and recovers near-optimal denominator bounds at every cell.

\bigskip

%-------------------------------------------------------------------------
\section{The 2nd-order Taylor pad in $w$}\label{sec:taylor-pad}

A chunk $Z$ of $w$-radius $r_Z$ admits a sharp interior midpoint $w_0$ at which
$\Heta, \HetaP, \HetaPP$ are evaluated by~\eqref{eq:price}. The chunk values are then
related to the midpoint values by 2nd-order Taylor expansion in $w$:
\begin{equation}\label{eq:taylor-pad}
\begin{aligned}
   \sup_{w\in Z}|\Phi(w)| &\le |\Phi(w_0)| + \tfrac{r_Z^2}{2}\,M_{\Phi''} + \tfrac{r_Z^3}{6}\,M_{\Phi'''}, \\
   \sup_{w\in Z}|\Phi'(w)| &\le |\Phi'(w_0)| + r_Z\,M_{\Phi''} + \tfrac{r_Z^2}{2}\,M_{\Phi'''}, \\
   \sup_{w\in Z}|\Phi''(w)| &\le \min\bigl(|\Phi''(w_0)| + r_Z\,M_{\Phi'''},\;M_{\Phi''}\bigr),
\end{aligned}
\end{equation}
where $M_{\Phi'''}$ is a Cauchy bound on the third derivative obtained from the
analyticity of $\Phi$ on a slightly larger disk. The padded values are sharper than
either the chunk-uniform bound $M_{\bullet}$ or the midpoint value alone, and are the
mechanism by which corner-subdivision rounds (cf.\ \S\ref{sec:subdivision}) drive the
certified sup down: each subdivision halves $r_Z$ and the pads shrink at order
$r_Z^2$, $r_Z$, and $1$ respectively, with the third estimate switching from the
Cauchy-padded value to the chunk-uniform bound at small $r_Z$.

\bigskip

%-------------------------------------------------------------------------
\section{Per-cell $\min(\text{interval},\text{norm})$}\label{sec:min-cell}

For cells where the merged numerator $2\Phi\Phi'D - DD\,\Phi^2$ has matched-magnitude
summands of opposite sign, the interval-arithmetic numerator
\(
   \mathrm{score}_{\mathrm{interval}}
   \;=\;
   \absub{2\Phi\Phi'D - DD\Phi^2}\big/\abslb{D}^3
\)
captures the destructive interference but suffers dependency-blow-up when the
constituents widen across the cell. The triangle-norm form
\(
   \mathrm{score}_{\mathrm{norm}}
   \;=\;
   2\,\absub{\Phi}\absub{\Phi'}/\abslb{D}^2 + \absub{DD}\absub{\Phi}^2/\abslb{D}^3
\)
discards interference but has no dependency-blow-up. Both quantities are valid upper
bounds on $|u_{pp}|$ on the cell. We always use the tighter of the two,
\(
   \mathrm{score}_{\mathrm{cell}}
   \;=\;
   \min(\mathrm{score}_{\mathrm{interval}},\ \mathrm{score}_{\mathrm{norm}}),
\)
which is again provably a valid upper bound. The trade-off shifts cell-by-cell:
near corner saturation, $\mathrm{score}_{\mathrm{norm}}$ wins; in the bulk,
$\mathrm{score}_{\mathrm{interval}}$ wins.

\bigskip

%-------------------------------------------------------------------------
\section{Symmetry of the V26 contour and corner subdivision}\label{sec:subdivision}

\paragraph{Four-fold symmetry.}
The function $\Heta$ at our parameters satisfies the two reflection identities
\(
   \Heta(\bar z) = \overline{\Heta(z)},
\)
\(
   \Heta(-z) = -\Heta(z).
\)
The first is a consequence of the real-coefficient Taylor expansion of $\Heta$ near
the origin; the second is the Hermite-physicist parity of the sign-pair $K_z$. The
joint four-fold symmetry reduces the chunk certifier's work by exactly~$4\times$:
only $w$-cells in the first quadrant of $\Om$ need to be evaluated primarily.

\paragraph{Corner subdivision.}
The chunk certifier uses a uniform $n_\bullet \times 4$ partition of $\dOm$, but the
sup is consistently dominated by chunks adjacent to the four corners
$\pm(1-\delta) \pm iB$. Recursively splitting each corner chunk $4$-way (or $2$-way for
1D edge chunks) and re-running the Taylor-pad evaluator on the sub-chunks halves
the certified sup roughly each round:
\[
\begin{array}{c|ccc}
   \text{round} & 1 & 2 & 3 \\\hline
   C_r & 2.27 & 1.23 & 0.383
\end{array}
\]
starting from $C_r=34.47$ at the unsubdivided 64-chunk baseline. The diminishing returns
reflect the fact that the binding cells move inwards from the corners as the local
bound becomes tighter than the bulk.

\bigskip

%-------------------------------------------------------------------------
\section{The $C_r$ landscape: formula $\times$ oracle constellation}\label{sec:landscape}

Each candidate $V_K$ in the refinement loop encodes a pair
$(\text{$C_r$ formula},\ \text{oracle constellation})$, not a linear sequence of
literals. After the truncation-free re-audit (\S\ref{sec:trunc-free}), the
\emph{patched} $C_r$ values rank as follows for our target parameters
$\eta=1/100$, $r=23/25$, $p_{\max}=1$:
\[
\begin{array}{c|l|c}
   V_K & C_r\ \text{formula} & C_r^{\mathrm{patched}} \\\hline
   V_{11} & C_r^{\phi,u_{zz}} = M_\phi M_{\phi'} I^2 + \tfrac12 \norm{u_{zz}} M_\phi^2 & 0.3255 \\
   V_{14} & C_r^{\phi,u_{zz},\,\mathrm{joint}} = J_1 + \tfrac12 \norm{u_{zz}} M_\phi^2 & \mathbf{0.2810} \\
   V_{15-24} & C_r^{\phi,u_{zz},\,\mathrm{joint2}} = J_1 + \tfrac12 J_2 & 0.4105 \\
   V_{25,26} & C_r^{\phi,u_{zz},\,\mathrm{merged}} = \tfrac12\,J_{\mathrm{merged}} & 0.3734
\end{array}
\]
($I=\sup 1/|D|$.)
The smallest $C_r^{\mathrm{patched}}$ is held by $V_{14}$. Two surprises emerge.

\paragraph{Joint-on-only-one-term.}
$V_{14}$'s mixed formula---$J_1$ (joint sup, term~1) plus $\norm{u_{zz}} M_\phi^2$
(triangle product, term~2)---wins over $V_{15\text{-}24}$'s
$J_1+\tfrac12 J_2$ (joint sup on \emph{both} terms). This is a consequence of the
form-A identity \eqref{eq:upp-merged}: replacing $\norm{u_{zz}} M_\phi^2$ by
$J_2 = \sup\norm{DD\,\Phi^2/D^3}$ changes which factor (the numerator $|DD|$ or the
correlation $|\Phi|^2$) is taken on its worst cell, and at our parameters the
correlation is more concentrated than the joint sup over a single cell can capture,
so the triangle bound on term~2 is paradoxically tighter on the patched oracles.

\paragraph{Merged loses (at finite resolution).}
$V_{25,26}$'s merged formula is exact at the level of \eqref{eq:upp-merged}, but the
certified sup $J_{\mathrm{merged}}$ at the chunk resolutions used for the audit
($n=2^{18}$) is still slightly above the triangle $2J_1+J_2$, by exactly the
crossover phenomenon in \S\ref{sec:formA-merged}. At higher resolution
($n\ge 2^{20}$ extrapolated), the merged would dominate; at the production resolution
of the audit, the joint mixture in $V_{14}$ wins.

\paragraph{Take-away.}
Beating any specific $C_r^{\mathrm{patched}}$ requires \emph{moving the audit floor},
not just tightening the literal:
\begin{enumerate}[leftmargin=2em]
   \item re-baseline the formula choice onto $V_{14}$ for an immediate $\sim 17\%$ win;
   \item subdivide chunks further (each round halved $J_{\mathrm{merged}}$ historically);
   \item replace the 2nd-order Taylor pad by a 4th-order pad with a tighter
         $M_{\Phi^{(4)}}$;
   \item use an alternative analyticity domain (octagonal, Cassini, \dots) that
         avoids corner saturation entirely.
\end{enumerate}
Items (3)--(4) are open analytic problems and require new oracles.

\bigskip

%-------------------------------------------------------------------------
\section{The Cauchy-radius lift and the analyticity-domain envelope}\label{sec:cauchy-lift}

The implicit-function family $\Finvmap$ is constructed in the original proof as the
fixed-point of a contraction on a Banach space of analytic functions on
$\mathrm{disk}(r_1\,r)$ with $r_1=(1+1/r)/2$. Our Lean axiomatisation supplies the
strictly weaker analyticity domain $\mathrm{disk}((1+r)/2)$ (which equals $0.96$ at
$r=23/25$), forcing the recentred Cauchy ball at any $z_0\in\mathrm{disk}\,r$ to have
radius $\rho\le(1-r)/2=0.04$. By the standard Cauchy estimate
\(
   |f^{(n)}(z_0)|
   \;\le\;
   n!\,\sup_{|z-z_0|=\rho} |f(z)|/\rho^n,
\)
shrinking $\rho$ inflates derivative bounds by $1/\rho^n$. Lifting the analyticity
domain to $\mathrm{disk}((3+r)/4) \approx 0.98$ enlarges the available Cauchy
radius to $\rho=(3-r)/4\approx 0.13$, a $1.5\times$ lift, with an $n$-th derivative
bound shrinking by $1.5^n$:
\begin{itemize}
   \item $1/D$ (first derivative): $1.5\times$ tighter,
   \item $u_{zz}$ via Cauchy on the second derivative: $2.25\times$ tighter,
   \item third-derivative Taylor pads (\S\ref{sec:taylor-pad}): $3.375\times$ tighter.
\end{itemize}
The wider domain claim $\mathrm{disk}((3+r)/4)$ is paper-citable in the original
proof's Step~4; bringing it into the Lean axiomatisation is the single largest
\emph{trust-expansion} gain available, with a downstream win factor that compounds
with the chunk-bound improvements above.

\bigskip

%-------------------------------------------------------------------------
\section{Summary of trace}

Across $34$ refinement sprints driven by the framework, the certified $C_r$ for
Lemma~5.4 traces (in pre-truncation-free units, then in truncation-free units after
the May-2026 patch):
\[
   8.2\!\times\!10^{14}
   \;\to\;
   1.13\!\times\!10^{9}
   \;\to\;
   3.79
   \;\to\;
   1
   \;\to\;
   1/268{,}000\quad
   \xrightarrow{\text{$K_{\max}$ patch}}\quad
   0.3255
   \;\to\;
   0.2857.
\]
The correspondence between the $C_r$ trace and the $\KG$-improvement trace
\eqref{eq:KGimp} is recorded in
\href{run:plots/cr_progression.pdf}{\texttt{plots/cr\_progression.pdf}} and
\href{run:plots/kg_improvement.pdf}{\texttt{plots/kg\_improvement.pdf}}. The non-
monotonic plateaus reflect formula switches; the steepest descents come from the joint
oracles (\S\ref{sec:joint-vs-triangle}) and the $K_{\max}$-truncation patch
(\S\ref{sec:trunc-free}).

The mathematical insights collected above---the auxiliary $z$-identity, the form-A
identity and merged-numerator interference, the joint-vs-triangle gap, the
$\eta^2$-tight $\Phi$-difference bound, the $\lambda$-decay parallelogram fix, the
denominator-fix and Taylor-pad machinery, the per-cell $\min(\text{interval},\text{norm})$,
the four-fold symmetry, and the analyticity-domain lift---are individually
modest, but their composition is what makes the certified $\KG$-improvement strictly
positive in arb-arithmetic at the $10^{-17}$ scale.

\end{document}
